\section{Introduction}
\label{sec:introduction}

For a von Neumann algebra $M\subseteq\B(H)$, write $M_1$ for its
closed operator-norm unit ball. The Kadison--Kastler distance between
two such algebras is
\begin{equation}\label{eq:kk-definition}
 d(M,N)=\max\left\{
 \sup_{x\in M_1}\inf_{y\in N_1}\norm{x-y},\quad
 \sup_{y\in N_1}\inf_{x\in M_1}\norm{x-y}\right\}.
\end{equation}
Kadison and Kastler introduced this metric in 1972
\cite{KadisonKastler1972}. Their perturbation problem asks whether
closeness of the unit balls forces spatial isomorphism, implemented by
a unitary close to the identity. We prove the uniform form of this
assertion.

\begin{theorem}[Strong Kadison--Kastler stability]\label{thm:main}
For every $\eps>0$ there is $\delta=\delta(\eps)>0$ such that the
following holds. If $H$ is any complex Hilbert space and
$M,N\subseteq\B(H)$ are unital von Neumann algebras with identity
$I_H$ and $d(M,N)<\delta$, there is a unitary $u\in\B(H)$ satisfying
\[
 uMu^*=N,\qquad \norm{u-I_H}<\eps.
\]
\end{theorem}

The tolerance is independent of the algebras, their type, their
representations, and the cardinalities of their preduals and Hilbert
spaces. Thus \Cref{thm:main} resolves the strong Kadison--Kastler
conjecture positively. Its hypothesis concerns the ordinary distance
\eqref{eq:kk-definition}.

\subsection{Historical context}

Early positive results treated type~I algebras:
Phillips proved spatial isomorphism, and Christensen obtained a conjugating
unitary close to the identity \cite{Phillips1974,Christensen1975}.
The work of Christensen, Johnson, and Raeburn--Taylor established the
amenable case \cite{Christensen1977,Johnson1977,RaeburnTaylor1977}.
Christensen subsequently developed the one-sided theory of near inclusions:
an amenable von Neumann algebra sufficiently nearly contained in another
can be moved into it by a small unitary \cite{Christensen1980}.
His earlier work also treated close von Neumann subalgebras of a common
finite von Neumann algebra \cite{Christensen1977II}. The common finite
ambient algebra is an essential hypothesis of that result; merely assuming
that both represented algebras are finite gives a different problem.
In a different common-ambient setting, Ino proved near-identity unitary
conjugacy for sufficiently close subalgebras of a countably decomposable
ambient algebra when both admit conditional expectations of finite
probabilistic index from that algebra \cite[Theorem~3.2]{Ino2016}.

The corresponding development for $C^*$-algebras clarifies the role of the
conclusion. Christensen, Sinclair, Smith, White, and Winter proved that
sufficiently close $C^*$-algebras acting on a separable Hilbert space
are spatially isomorphic if one of them is norm separable and nuclear
\cite{CSSWW2010,CSSWW2012}.
Johnson's examples exclude a uniform bound on the displacement of an
implementing unitary that tends to zero with the algebra distance;
they do not exclude spatial isomorphism \cite{Johnson1982}.

For nonamenable von Neumann algebras, Cameron, Christensen, Sinclair, Smith,
White, and Wiggins constructed strongly stable factors of the form
$(L^\infty(X)\rtimes\mathrm{SL}_n(\mathbb Z))\barotimes R$, where
$n\geq3$, the action is free, ergodic, and probability preserving on a
standard nonatomic probability space, and $R$ is the hyperfinite
type~$\mathrm{II}_1$ factor \cite{CCSSWW2012,CCSSWW2014}.
Their proof combines amenable near-inclusion methods, a common standard
representation, and bounded cohomology. Their later work transfers
property~$\Gamma$ and Cartan structure between sufficiently close type
$\II_1$ factors with separable predual, and also studies stability of
solidity, strong solidity, and tensor-product structure
\cite{CCSSWW2017}.

The passage between representations is tied to Kadison's similarity
problem: whether a bounded homomorphism from a $C^*$-algebra into
$\B(H)$ is conjugate by an invertible operator to a $*$-homomorphism.
Christensen's work on property~$\Gamma$ factors provides a
foundational special case \cite{Christensen1986,Christensen2001}.
Christensen, Sinclair, Smith, and White established perturbation results
for the similarity property and operator-algebraic invariants
\cite{CSSW2010}; Cameron and collaborators then proved that the similarity
property is equivalent to continuity of commutants, with the required
uniformity over representations \cite{CCSSWW2013}.
These results explain why ordinary norm closeness, control at all matrix
levels, and implementation on the original Hilbert space must be kept
distinct throughout a proof of strong stability.
The companion manuscript \cite{OpenAI2026Similarity} also proves a
uniform commutator estimate, using invariant tensors and a free-corner
Hankel test. The amplification argument here instead uses free-word
randomization and is proved independently. Neither commutator estimate
alone constructs a nearby multiplicative correspondence between two
algebras; Appendix~\ref{app:comparison} compares their scope and this
remaining perturbation step.

The unit-ball metric supplies nearby elements but does not specify a
linear or multiplicative correspondence. We must construct such a
correspondence and control its implementation on the given Hilbert space,
with estimates independent of the algebras and representations.

\subsection{Proof overview}

We first work on separable Hilbert spaces. For finite factors, two
arguments meet in \Cref{sec:finite}. The first supplies matrix norm
control: for a separably acting type $\II_1$ factor $M\subseteq\B(H)$
and $h=h^*$, \Cref{thm:amplification} bounds all matrix amplifications
of $x\mapsto[h,x]$ by an absolute multiple of its ordinary operator
norm. It combines Haagerup's cyclic similarity estimate
\cite{Haagerup1983} and Popa's independence theorem for tracial
ultraproducts \cite{Popa2014} with a free-word randomization argument.
The constant is independent of matrix size and representation.

A separate argument constructs an isomorphism close in ordinary norm
to the inclusion. Using the common-position construction of
\cite[Lemma~4.10]{CCSSWW2014}, we move the factors to a tracial
representation in which one vector represents both traces, and the
left and right actions of a common hyperfinite subfactor $R$ agree.
The inclusion is irreducible, meaning $R'\cap M=\C I$.

For a countable strongly dense subgroup $G$ of the unitary group of $M$,
choose nearby unitaries $v_g\in N$. Their multiplication errors are
uniformly small. Section~\ref{sec:boundary} constructs a probability
space $D$ with an amenable $G$-action preserving null sets. An adaptation
of the averaging and polar-correction method for approximate
representations \cite{Kazhdan1982,BurgerOzawaThom2013} corrects the
nearby unitaries to measurable fields $W_g:D\to\U(N)$ satisfying
\[
 W_g(x)W_h(g^{-1}x)=W_{gh}(x).
\]
Here $G$ need not be amenable: the correction produces a cocycle over
its amenable action, not a representation on the original space.
On $D^2$, a tracial averaging argument and double ergodicity force a
suitable equivariant vector field to be constant. We give the
unitary-coefficient double-ergodicity argument in full, using the
bilateral random-walk viewpoint of Kaimanovich; Bj\"orklund gives
the corresponding unitary formulation and another proof
\cite{Kaimanovich2003,Bjorklund2014}. The resulting linear bijection preserves
adjoints and squares and is close to the inclusion; proximity excludes
reversal of multiplication. Transporting this isomorphism back and
applying the commutator estimate gives a small unitary on the original
Hilbert space.

For infinite factors, a common copy of $\B(\ell^2)$ identifies the
ordinary gap with the supremum of all matrix gaps. Tensoring with one
fixed hyperfinite factor gives type $\III_1$ factors with an
irreducible hyperfinite finite subfactor admitting a conditional
expectation. This preparation uses Marrakchi's absorption criterion
for a trivial bicentralizer and Haagerup's construction of the expected
hyperfinite subfactor \cite{Marrakchi2020,Haagerup1987}.
Houdayer and Marrakchi's September 2026 preprint proves triviality of
the bicentralizer for every type $\III_1$ factor with a faithful normal
state \cite[Theorem~G]{HoudayerMarrakchi2026}; the preparation here uses
only the earlier absorption criterion.

Small adjustments align the subfactor's left and right actions and
provide a common cyclic separating vector $\Omega$ whose two vector
states centralize $R$. For either algebra, the closure of
$x\Omega\mapsto x^*\Omega$ is its Tomita operator.
\Cref{prop:modular-graphs} shows that small matrix gaps force its two
graphs to be close when the output is multiplied by any $c>0$,
uniformly in $c$. This uniformity permits comparison of the modular
actions. A uniform infinitesimal estimate for type $\III$ factors
controls norm-continuous paths of unitaries whose conjugations nearly
preserve the algebra. Using Arveson's nest distance formula
\cite{Arveson1975}, analytic smoothing brings the modular comparison
into this setting. Limits of conjugation maps yield, at each time, a
unitary whose conjugation preserves one algebra and which is close to
the other algebra's modular unitary. Averaging corrects their
multiplication law. Section~\ref{sec:modular-alignment} thereby obtains
a single small change of position making the first algebra invariant
under that modular group.

We form crossed products using this common action. Their unit balls
are close at every matrix size, and continuous decomposition identifies
one crossed product as a type $\II_\infty$ factor
\cite{Takesaki1973,Takesaki2003II}. After aligning matrix units, its finite corner
and the corresponding nearby corner are both finite factors. The
finite-factor theorem conjugates these corners and hence the crossed
products. Averaging gives a unital completely positive (UCP) map between
the base algebras, meaning that it preserves the identity and positivity
at every matrix size. A state slice gives a second UCP map between the
original factors. The criterion in \Cref{prop:ucp}, derived from
Ricard--Roydor's multiplication perturbation theorem
\cite[Proposition~3.2]{RicardRoydor2014}, turns both maps into small
spatial isomorphisms. It does not require the averaging maps to be normal.

Finally, aligning centers and selecting measurable unitaries assembles
the factor results on separable spaces. For an arbitrary Hilbert space,
finite sets of operators and vectors lie in suitable separable
restrictions. Extend their small conjugating unitaries by the identity.
A subnet of the conjugation maps converges pointwise in the weak
operator topology to a UCP map. The finite tests put its range in the
required algebra, and the same criterion finishes the proof.
All tolerances depend only on the requested final unitary displacement.

\Cref{fig:proof} displays how the finite-factor result enters both the
finite and infinite cases and how they meet in the final reductions.

\begin{figure}[htbp]
\centering
\begin{tikzpicture}[
 box/.style={draw,rounded corners=2pt,align=center,text width=3.55cm,
             minimum height=1.0cm,inner sep=5pt,font=\small},
 arr/.style={-{Stealth[length=2mm]},thick}]
 \node[box] (amp) at (-2.5,0) {Commutator amplification\\Section~\ref{sec:amplification}};
 \node[box] (boundary) at (2.5,0) {Boundary and common trace\\Sections~\ref{sec:boundary}--\ref{sec:finite}};
 \node[box] (finite) at (0,-1.65) {Finite factors on separable spaces\\Theorem~\ref{thm:finite}};
 \node[box] (modular) at (-4,-3.8) {Preparation and modular alignment\\Sections~\ref{sec:modular-preparation}--\ref{sec:modular-alignment}};
 \node[box] (corner) at (1,-3.8) {Finite crossed-product corner\\and two descents\\Section~\ref{sec:cores}};
 \node[box] (factors) at (1,-5.8) {All separably acting factors\\Theorem~\ref{thm:factors}};
 \node[box] (center) at (-2.5,-7.6) {Common center and\\measurable assembly\\Section~\ref{sec:reductions}};
 \node[box] (arbitrary) at (2.5,-7.6) {Finite tests on arbitrary $H$\\Section~\ref{sec:reductions}};
 \draw[arr] (amp)--(finite);
 \draw[arr] (boundary)--(finite);
 \draw[arr] (finite)--(corner);
 \draw[arr] (modular)--(corner);
 \draw[arr] (corner)--(factors);
 \draw[arr] (finite.east) -- ++(2.0,0) |- (factors.east);
 \draw[arr] (factors)--(center);
 \draw[arr] (center)--(arbitrary);
\end{tikzpicture}
\caption{The finite-factor theorem has two inputs and is reused in the
crossed-product argument. The UCP criterion supplies pointwise
normalizers in Section~\ref{sec:modular-alignment}, both descents in
Section~\ref{sec:cores}, and the final limit in Section~\ref{sec:reductions}.}
\label{fig:proof}
\end{figure}

Section~\ref{sec:foundations} proves the transfer tools.
Sections~\ref{sec:amplification}--\ref{sec:finite} establish the finite
case. Sections~\ref{sec:modular-preparation}--\ref{sec:cores} treat
infinite factors. Section~\ref{sec:reductions} removes the factor and
separability assumptions and gathers the final quantifiers.
Appendix~\ref{app:comparison} distinguishes uniformity over
representations from uniformity over factors and compares the
commutator methods of this paper and the similarity companion.
